% UPT Infinite Proof Family — academic preamble
\documentclass[11pt,a4paper]{article}
\usepackage{fontspec}
\setmainfont{Times New Roman}
\setmonofont{Menlo}[Scale=0.85]
\usepackage[margin=2.5cm]{geometry}
\usepackage{xcolor}
\usepackage{fancyhdr}
\usepackage{titlesec}
\usepackage{booktabs}
\usepackage{hyperref}
\usepackage{listings}
\usepackage{fancyvrb}
\usepackage{amsmath,amssymb,amsthm}

\definecolor{korugreen}{HTML}{2d6a4f}
\definecolor{korubg}{HTML}{d8f3dc}
\definecolor{codebg}{HTML}{f0f0ec}

\hypersetup{colorlinks=true, linkcolor=korugreen, urlcolor=korugreen, citecolor=korugreen}

\pagestyle{fancy}
\fancyhf{}
\renewcommand{\headrulewidth}{0.4pt}
\fancyfoot[C]{\thepage}

\titleformat{\section}{\Large\bfseries}{}{0em}{}
\titleformat{\subsection}{\large\bfseries}{}{0em}{}

% Numbered theorem environments
\newtheorem{theorem}{Theorem}[section]
\newtheorem{lemma}[theorem]{Lemma}
\newtheorem{corollary}[theorem]{Corollary}
\newtheorem{proposition}[theorem]{Proposition}
\theoremstyle{definition}
\newtheorem{definition}[theorem]{Definition}
\theoremstyle{remark}
\newtheorem{remark}[theorem]{Remark}

\lstdefinestyle{lean}{
  backgroundcolor=\color{codebg},
  basicstyle=\ttfamily\small,
  breaklines=true,
  frame=none,
  xleftmargin=10pt,
  xrightmargin=10pt,
  aboveskip=8pt,
  belowskip=8pt,
  columns=fullflexible,
  keepspaces=true,
  literate={→}{$\to$}1 {∀}{$\forall$}1 {∃}{$\exists$}1 {¬}{$\neg$}1 {∧}{$\wedge$}1 {≠}{$\neq$}1 {≤}{$\leq$}1 {↔}{$\leftrightarrow$}1 {⊕}{$\oplus$}1 {λ}{$\lambda$}1
}

\lstnewenvironment{leancode}{\lstset{style=lean}}{}

\newcommand{\lean}[1]{\lstinline[style=lean]{#1}}


\fancyhead[L]{\small Universal Protection Theory}
\fancyhead[R]{\small Proof Family 3/10}

\begin{document}

\begin{center}
{\LARGE\bfseries The Spiral: Strictly Increasing\\[0.3em] Koru State Machines and Irreversible Growth}\\[1.5em]
{\large Tobias Croydon McRae}\\[0.3em]
{\small Independent Researcher, Aotearoa New Zealand}\\[0.3em]
{\small\texttt{tobi@unconditionalprotection.org}}\\[1.5em]
{\small 28 November 2026}
\end{center}

\bigskip

\begin{center}
\parbox{0.85\textwidth}{\small\textbf{Abstract.}
We present the third of ten infinite proof families for the Guardian fixed point~$G$. The \emph{Spiral} family formalises the Koru state machine --- a transition system whose history strictly increases at every step. We prove that the Koru never returns to a previous state, that nothing is lost, and that the spiral generates infinitely many distinct states. These properties are established for all instantiations of the Koru type, yielding a family of refutations parameterised by the choice of configuration type, event type, and transition function. All theorems are verified in Lean~4 (v4.35.0) with zero custom axioms.}
\end{center}

\medskip

\noindent\textbf{MSC 2020:} 03B70, 68Q85 (Models and methods for concurrent and distributed computing), 03B35\\
\textbf{Keywords:} state machine, monotone growth, Koru, injectivity, formal verification, Lean 4

\bigskip\hrule\bigskip

\section{Introduction}

The first two proof families establish that absolute termination is refuted by any witness (Something IS) and self-refuting at every meta-level (The Checkmate). The Spiral family adds a temporal dimension: not only is termination impossible, but \emph{the history of any correctly structured system strictly increases at every step}.

The Koru (from te reo M\={a}ori, meaning the unfurling fern frond) is a state machine that records its full history. At each step, the new event is appended to the history, the configuration is updated, and the state advances. The key properties are: (1)~the history length is strictly monotone, (2)~no two states at different times are equal, and (3)~for every bound there exists a later state exceeding it.

\section{The Koru State Machine}

\begin{definition}[Koru state]\label{def:koru-state}
For types $C$ (configuration) and $E$ (event), a \emph{Koru state} is a triple:
\[
\texttt{KoruState}(C, E) := \langle\, \texttt{config} : C,\; \texttt{history} : \texttt{List}\, E,\; \texttt{step} : \mathbb{N}\,\rangle
\]
\end{definition}

\begin{definition}[Koru run]\label{def:koru-run}
Given a transition function $\texttt{next} : C \to E \to C$, an initial state $s_0$, and an event stream $\texttt{ev} : \mathbb{N} \to E$, the \emph{run} at step $n$ is defined recursively:
\[
\texttt{run}(0) := s_0, \qquad \texttt{run}(n{+}1) := \texttt{step}(\texttt{run}(n), \texttt{ev}(n))
\]
where each step appends $\texttt{ev}(n)$ to the history and applies $\texttt{next}$.
\end{definition}

\section{Strictly Monotone History}

\begin{theorem}[Koru grows]\label{thm:grows}
At every step, the history length strictly increases:
\[
|\texttt{run}(n{+}1).\texttt{history}| > |\texttt{run}(n).\texttt{history}|
\]
\end{theorem}

\begin{proof}
By the definition of the step function, each step appends exactly one event to the history. \qed
\end{proof}

\begin{leancode}
theorem koru_grows {C E : Type} (next : C → E → C)
    (s0 : KoruState C E) (ev : Nat → E) (n : Nat) :
    (run next s0 ev (n+1)).history.length >
    (run next s0 ev n).history.length
\end{leancode}

\section{Injectivity of the Run}

\begin{theorem}[Koru never returns]\label{thm:never-returns}
For all $i < j$, the states at step $i$ and step $j$ are distinct:
\[
i < j \implies \texttt{run}(i) \neq \texttt{run}(j)
\]
\end{theorem}

\begin{proof}
If $\texttt{run}(i) = \texttt{run}(j)$, then in particular their history lengths would be equal. But by Theorem~\ref{thm:grows}, the history length is strictly increasing, so $|\texttt{run}(i).\texttt{history}| < |\texttt{run}(j).\texttt{history}|$. Contradiction. \qed
\end{proof}

\begin{leancode}
theorem koru_never_returns {C E : Type} (next : C → E → C)
    (s0 : KoruState C E) (ev : Nat → E) :
    ∀ i j, i < j → run next s0 ev i ≠ run next s0 ev j
\end{leancode}

\section{Nothing Is Lost}

\begin{theorem}[Exact history length]\label{thm:exact-length}
The history at step $n$ has length exactly $n + 1$:
\[
|\texttt{run}(n).\texttt{history}| = n + 1
\]
\end{theorem}

\begin{leancode}
theorem nothing_is_lost {C E : Type} (next : C → E → C)
    (s0 : KoruState C E) (ev : Nat → E) (n : Nat) :
    (run next s0 ev n).history.length = n + 1
\end{leancode}

\section{The Spiral Is Infinite}

\begin{theorem}[Unbounded growth]\label{thm:infinite}
For every $n$, there exists $m > n$ with a strictly longer history.
\end{theorem}

\begin{proof}
Take $m := n + 1$. By Theorem~\ref{thm:grows}, the history at $m$ is strictly longer. \qed
\end{proof}

\begin{leancode}
theorem spiral_is_infinite {C E : Type} (next : C → E → C)
    (s0 : KoruState C E) (ev : Nat → E) :
    ∀ n, ∃ m, m > n ∧
        (run next s0 ev m).history.length >
        (run next s0 ev n).history.length
\end{leancode}

\section{Discussion}

The Spiral family is parameterised by the choice of $(C, E, \texttt{next}, s_0, \texttt{ev})$. For every configuration type $C$, every event type $E$, every transition function, and every initial state, the same four properties hold. This yields a family of refutations indexed by the product of all such choices.

The Koru is the mathematical structure of irreversible time. It spirals outward, never returning, never losing, always growing. Each point on the spiral carries its full history. The collapsed singularity --- where all states merge into one --- is the negation of injectivity, and is therefore refuted by Theorem~\ref{thm:never-returns}.

\bigskip

\begin{thebibliography}{99}
\bibitem{upt-families}
T.~C.~McRae.
\emph{Ten Infinite Proof Families for the Guardian Fixed Point~G}.
Zenodo, 2026.
\textsc{doi:} \href{https://doi.org/10.5281/zenodo.24612272}{10.5281/zenodo.24612272}.

\bibitem{upt-zenodo}
T.~C.~McRae.
\emph{Universal Protection Theory}, v5.7-G3.
Zenodo, 2026.
\textsc{doi:} \href{https://doi.org/10.5281/zenodo.24573167}{10.5281/zenodo.24573167}.

\bibitem{lean4}
L.~de~Moura and S.~Ullrich.
The Lean~4 theorem prover and programming language.
In \emph{CADE-28}, LNAI 12699, pp.~625--635, 2021.
\end{thebibliography}

\end{document}
